% --- Avant-propos -----------------------------------------------------------
\chapter*{Avant-propos}
\markboth{Avant-propos}{Avant-propos}
\addcontentsline{toc}{chapter}{Avant-propos}

Ce livre est le deuxième de la collection Boky Matematika, et le premier
destiné à la série OSE. Il reprend, sans en changer un principe, ce qui a
fait le livre de Terminale S : un même texte pour l'élève et pour le
professeur, un cours qui ne mobilise jamais une notion avant qu'elle ait été
enseignée, et une banque d'exercices exclusivement tirée de sujets réellement
posés.

La série OSE — Organisation Société Économie — n'est pas une version allégée
de la série S. Elle a son propre équilibre : l'analyse y occupe une place
plus large encore (quatre-vingts heures), l'arithmétique et le calcul
matriciel en disparaissent complètement au profit d'un unique système linéaire
à trois inconnues, et les mathématiques financières — absentes du programme
de la série S, et absentes elles-mêmes du tableau des quatre composantes
chiffrées du Programme d'Études — trouvent leur place dans une annexe dédiée
plutôt que parmi les parties numérotées : c'est ainsi qu'elles figurent dans
la Répartition Annuelle et dans chaque sujet de baccalauréat examiné, sans
jamais apparaître comme une composante à part entière du programme officiel.
La partie Géométrie, elle, surprend par sa richesse : le Programme d'Études
2026 y place les nombres complexes, le barycentre et les similitudes, un
contenu presque aussi dense que celui de la série S, alors même qu'aucun
sujet de baccalauréat OSE examiné à ce jour ne l'a encore interrogé. Ce livre
choisit d'y rester fidèle : le programme officiel fait foi, et les élèves qui
l'auront travaillé ne seront jamais pris au dépourvu si les sujets s'alignent
un jour sur lui.

\section*{Un mot à l'élève de terminale}

Vous entrez dans l'année où les mathématiques cessent d'être une succession
de techniques pour devenir un outil de décision. Une fonction de coût, un
capital placé, une série statistique à deux variables : ce sont les mêmes
raisonnements qu'en analyse pure, appliqués à des questions que vous
croiserez hors de la salle d'examen.

Vous remarquerez que chaque chapitre commence par une page de révision qui ne
parle jamais de la notion nouvelle. C'est voulu. On ne vous demandera jamais
de deviner ce que vous n'avez pas encore appris : on vérifie d'abord que les
outils de l'an dernier répondent encore, puis on ouvre le cours. Prenez ces
pages au sérieux, elles coûtent dix minutes et vous en font gagner beaucoup.

\section*{Un mot au professeur}

Les volumes horaires indiqués en tête de chaque partie sont ceux du Programme
d'Études — quatre-vingts heures pour l'analyse, vingt-cinq pour l'algèbre,
quarante-cinq pour la géométrie, trente pour le traitement des données. Les
mathématiques financières n'y figurent pas en tant que composante ; c'est
pourquoi ce livre les présente en annexe plutôt que comme une cinquième
partie numérotée, bien que leur volume horaire ne soit fixé nulle part dans
le PE, que la Répartition Annuelle les place en quatrième période, et que
chaque sujet de baccalauréat OSE examiné leur consacre un exercice entier,
coefficient compris. Le tableau de progression annuelle, en annexe
professeur, propose un calendrier compatible avec les cinq périodes du
calendrier scolaire, annexe financière comprise.

Aucun sujet de concours ENI ou ESPA ne figure dans ce livre : ces concours
recrutent sur un profil scientifique qui ne correspond pas à la série OSE.
La banque d'annales s'appuie sur les sessions malgaches conservées (2021,
2022 et 2026) et, ponctuellement, sur des séries économiques francophones
(Côte d'Ivoire G1/G2, Sénégal G) lorsque la banque malgache seule ne suffit
pas à couvrir un chapitre.

Les démonstrations sont toutes écrites. Cela ne veut pas dire qu'elles
doivent toutes être faites au tableau : à vous de choisir. Elles sont là pour
que l'élève qui veut comprendre pourquoi un résultat est vrai trouve la
réponse dans son livre plutôt que nulle part.

\vspace{1.2em}
\noindent\textit{Antananarivo}
